\section{Drøfting}
\label{sec:discussion-nn}

\subsection{Kva teorien forklarer som forgjengarane ikkje gjer}

Teorien gjev sameina forklaringar på fenomen som eksisterande rammeverk adresserer berre delvis eller ikkje i det heile.

\textbf{Kvifor stil predikerer meir enn materiale.} Tilpassingslandskapet forklarer det empiriske funnet direkte: stil er ein stedfortreder for det fulle trykk-settet (minus den klassekonstante affordansen), medan materiale er eitt trykk. Den aggregerte variabelen ber naudsynleg meir informasjon om formvariasjon enn nokon einskildkomponent.

\textbf{Kvifor stilar er gradientar, ikkje klynger.} Hadde stilar vore pålagde av eksterne kategoriske krefter, ville dei produsert diskrete klynger i morforommet. Dei negative silhouette-skårane ($-0{,}338$) stadfestar at stilar er gradientfenomen: kontinuerlege overlappande fordelingar produserte av agentar som navigerer eit delt men skiftande landskap. Stil er det statistiske sporet av navigasjonen, ikkje ei årsak til han \parencite{Kubler1962}.

\textbf{Kvifor formhistoria ligg etter materialhistoria.} Stål fanst i hundreår før røyrstolen. Teorien forklarer denne latensen gjennom skiljet mellom materialaffordanse (operasjonane eit materiale \textit{tillèt}) og sjølve operasjonen (teknikken som \textit{aktiverer} affordansen). Materialtilgang er naudsynleg men ikkje tilstrekkeleg; operasjonen må tre inn i kunnskapsrommet $K$ før morforommet kan ekspandere inn i den nyleg tilgjengelege regionen.

\textbf{Kvifor «form følgjer funksjon» er empirisk usant innanfor ein klasse.} Affordansestrukturen definerer klassen og er konstant innanfor han (proposisjon 1). Per definisjon ber han null informasjon om variasjon innanfor klassen. All variasjon må vere driven av trykk utover den klassedefinerande affordansen.

\subsection{Avgrensingar}

Teorien er testa mot éin einaste artefaktklasse (den europeiske stolen). Generalisering til andre klassar; arkitektur, grafisk design, programvaregrensesnitt, biologisk design; er teoretisert men ikkje enno demonstrert.

Korpuset er geografisk avgrensa til europeiske samlingar. Ikkje-europeiske designtradisjonar er ikkje representerte. Dei tredimensjonale mesh-modellane er utrekna frå katalogbilete, ikkje direkte skanna frå fysiske objekt. Direkte 3D-skanning ville styrka det empiriske grunnlaget.

Teorien spesifiserer ikkje enno den funksjonelle forma til tilpassingslandskapet. Vektene $w_i(t)$ i landskapslikninga er behandla som tidsvarierande parametrar, ikkje som funksjonar med spesifisert dynamikk.

\subsection{Tilhøve til eksisterande designteoriar}

Teorien er \textit{sterkare} enn C-K-teori åleine: han legg morforomsgeometri og landskapsdynamikk til den epistemiske logikken. Han er \textit{meir generell} enn formgrammatikkar åleine: han legg agenten, seleksjonstrykket og landskapet til det syntaktiske apparatet. Han er \textit{meir spesifikk} enn Levin sitt multiskala-kompetanserammeverk åleine: han gjev ein designdomene-formalisering som gjer det generelle rammeverket empirisk operasjonelt for artefaktar.

Teorien \textit{erstattar} ikkje desse rammeverka; han \textit{integrerer} dei. Kjernalikninga $\mathrm{Form} = SG(\nabla_{C(A)} L(c,t))$ er integrasjonspunktet.
